% Part: first-order-logic
% Chapter: syntax-and-semantics
% Section: assignments

\documentclass[../../../include/open-logic-section]{subfiles}

\begin{document}

\olfileid{fol}{syn}{ass}

\olsection{ચલ-નિયુક્તિઓ}

\begin{explain}
!!{variable} નિયુક્તિ~$s$ \emph{દરેક} ચલને મૂલ્ય આપે છે---અને ચલો
અનંત સંખ્યામાં છે. અલબત્ત, આ જરૂરી નથી. આપણે !!{variable}
નિયુક્તિઓને બધાં !!{variable}s માટે મૂલ્યો નિયુક્ત કરવાનું માત્ર એટલા
માટે કહીએ છીએ કે તેનાથી બાબતો ઘણી સરળ બને છે. પદ~$t$નું મૂલ્ય અને
!!a{formula}~$!A$,~!!a{structure}માં~$s$ સાપેક્ષ સંતોષાય છે કે નહિ,
માત્ર~$t$માંના !!{variable}sને અને~$!A$ના મુક્ત !!{variable}sને~$s$
કરેલી નિયુક્તિઓ પર આધાર રાખે છે. આગળનાં બે વિધાનો આ જ કહે છે.
``આધાર રાખે છે'' એ વિચાર ચોક્કસ કરવા આપણે બતાવીએ છીએ કે~$t$માંના બધા
ચલો પર સંમત થતી કોઈ પણ બે ચલ-નિયુક્તિઓ સરખું મૂલ્ય આપે છે; અને જો બે
ચલ-નિયુક્તિઓ~$!A$ના બધા મુક્ત ચલો પર સંમત હોય, તો~$!A$ એક નિયુક્તિ
સાપેક્ષ સંતોષાય ત્યારે અને માત્ર ત્યારે બીજી નિયુક્તિ સાપેક્ષ સંતોષાય.
\end{explain}

\begin{prop}\ollabel{prop:valindep}
જો પદ~$t$માંના !!{variable}s $x_1$, \dots,~$x_n$માંના હોય અને
$i=1$, \dots,~$n$ માટે $s_1(x_i)=s_2(x_i)$ હોય, તો
$\Value{t}{M}[s_1]=\Value{t}{M}[s_2]$.
\end{prop}

\begin{proof}
~$t$ની સંકુલતા પર અનુમાન વડે. પાયાના કિસ્સામાં~$t$ !!a{constant}
અથવા $x_1$, \dots,~$x_n$માંથી કોઈ એક ચલ હોઈ શકે. જો $t=c$ હોય, તો
$\Value{t}{M}[s_1]=\Assign{c}{M}=\Value{t}{M}[s_2]$. જો $t=x_i$ હોય,
તો વિધાનની પૂર્વધારણાથી $s_1(x_i)=s_2(x_i)$, અને તેથી
$\Value{t}{M}[s_1]=s_1(x_i)=s_2(x_i)=\Value{t}{M}[s_2]$.

અનુમાનાત્મક પગલા માટે ધારો કે $t=\Atom{f}{t_1,\dots,t_k}$ અને દાવો
$t_1$, \dots,~$t_k$ માટે સાચો છે. ત્યારે
\begin{align*}
  \Value{t}{M}[s_1] & = \Value{\Atom{f}{t_1, \dots, t_k}}{M}[s_1] \\
  & = \Assign{f}{M}(\Value{t_1}{M}[s_1], \dots, \Value{t_k}{M}[s_1]).
\intertext{$j=1$, \dots,~$k$ માટે~$t_j$ના !!{variable}s
    $x_1$, \dots,~$x_n$માંના છે. અનુમાનાત્મક પૂર્વધારણાથી
    $\Value{t_j}{M}[s_1]=\Value{t_j}{M}[s_2]$. તેથી,}
\Value{t}{M}[s_1] & = \Value{\Atom{f}{t_1, \dots, t_k}}{M}[s_1] \\
  & = \Assign{f}{M}(\Value{t_1}{M}[s_1], \dots, \Value{t_k}{M}[s_1]) \\
  & = \Assign{f}{M}(\Value{t_1}{M}[s_2], \dots, \Value{t_k}{M}[s_2]) \\
  & = \Value{\Atom{f}{t_1, \dots, t_k}}{M}[s_2] = \Value{t}{M}[s_2].
\end{align*}
\end{proof}

\begin{prop}\ollabel{prop:satindep}
જો~$!A$ના મુક્ત !!{variable}s $x_1$, \dots,~$x_n$માંના હોય અને
$i=1$, \dots,~$n$ માટે $s_1(x_i)=s_2(x_i)$ હોય, તો
$\Sat{M}{!A}[s_1]$ ત્યારે અને માત્ર ત્યારે જ્યારે
$\Sat{M}{!A}[s_2]$.
\end{prop}

\begin{proof}
આપણે~$!A$ની સંકુલતા પર અનુમાન વાપરીએ છીએ. પાયાના કિસ્સામાં, જ્યાં
$!A$ આણ્વિક છે, ત્યાં~$!A$ આમાંથી કોઈ હોઈ શકે:
\iftag{prvTrue}{$\ltrue$,}{}
\iftag{prvFalse}{$\lfalse$,}{}
$k$-સ્થાની વિધેય~$R$ અને પદો $t_1$, \dots,~$t_k$ માટે
$\Atom{R}{t_1,\dots,t_k}$, અથવા પદો $t_1$ અને~$t_2$ માટે
$\eq[t_1][t_2]$. છેલ્લા બે કિસ્સામાં આપણે !!{biconditional}ની માત્ર આગળની દિશા
બતાવીએ છીએ, કારણ કે ઊલટી દિશાની સાબિતી સમમિત છે.

\begin{enumerate}
\tagitem{prvTrue}{%
  \indcase{!A}{\ltrue}{$\Sat{M}{\indfrm}[s_1]$ અને
    $\Sat{M}{\indfrm}[s_2]$ બંને.}}{}

\tagitem{prvFalse}{%
  \indcase{!A}{\lfalse}{$\Sat/{M}{!A}[s_1]$ અને
    $\Sat/{M}{!A}[s_2]$ બંને.}}{}

\item
  \indcase{!A}{\Atom{R}{t_1, \ldots, t_k}}{ધારો કે
    $\Sat{M}{\indfrm}[s_1]$. ત્યારે
    \[
    \langle \Value{t_1}{M}[s_1], \ldots, \Value{t_k}{M}[s_1] \rangle
    \in \Assign{R}{M}.
    \]
    $i=1$, \dots,~$k$ માટે \olref{prop:valindep}થી
    $\Value{t_i}{M}[s_1]=\Value{t_i}{M}[s_2]$. તેથી આપણને
    $\langle \Value{t_1}{M}[s_2], \ldots, \Value{t_k}{M}[s_2] \rangle
    \in \Assign{R}{M}$ પણ મળે છે, અને તેથી
    $\Sat{M}{\indfrm}[s_2]$.}
\sourcecorrection{OLSEM-009}{મૂળની \texttt{assignments.tex},
પંક્તિ~91માં અંતિમ યુગ્મનું પ્રથમ સ્થાન~$t_i$ છે, જેથી એક મુક્ત
સૂચકચલ રહી જાય છે અને યુગ્મના પહેલા પદથી શરૂ થતી યાદી મળતી નથી.
ઘટકવાર સમાનતા લાગુ કર્યા પછીનું યુગ્મ $t_1$થી~$t_k$ સુધીનું જ છે;
તેથી પહેલા સ્થાનમાં~$t_1$ પુનઃસ્થાપિત કર્યું છે.}
\item
  \indcase{!A}{\eq[t_1][t_2]}{ધારો કે $\Sat{M}{\indfrm}[s_1]$.
    ત્યારે $\Value{t_1}{M}[s_1]=\Value{t_2}{M}[s_1]$. તેથી,
    \begin{align*}
      \Value{t_1}{M}[s_2] & = \Value{t_1}{M}[s_1]
      & \text{(\olref{prop:valindep}થી)} \\
      & = \Value{t_2}{M}[s_1]
      & \text{($\Sat{M}{\eq[t_1][t_2]}[s_1]$ હોવાથી)}\\
      &= \Value{t_2}{M}[s_2]
      & \text{(\olref{prop:valindep}થી),}
    \end{align*}
    તેથી $\Sat{M}{\eq[t_1][t_2]}[s_2]$.}
\end{enumerate}

હવે~$!A$ કરતાં ઓછી સંકુલતા ધરાવતા બધા !!{formula}s~$!B$ માટે
$\Sat{M}{!B}[s_1]$ ત્યારે અને માત્ર ત્યારે જ્યારે $\Sat{M}{!B}[s_2]$
એમ ધારો. અનુમાનાત્મક પગલું~$!A$ના મુખ્ય સંક્રિયકથી નક્કી થતા કિસ્સાઓ
પ્રમાણે આગળ વધે છે. દરેક કિસ્સામાં આપણે !!{biconditional}ની માત્ર આગળની દિશા
બતાવીએ છીએ; ઊલટી દિશાની સાબિતી સમમિત છે. પરિમાણકોના કિસ્સા સિવાયના
બધા કિસ્સામાં આપણે~$!A$નાં ઉપ-!!{formula}s~$!B$ને અનુમાનાત્મક
પૂર્વધારણા લાગુ કરીએ છીએ.~$!B$ના મુક્ત ચલો~$!A$ના મુક્ત ચલોમાંના છે.
તેથી જો $s_1$ અને~$s_2$,~$!A$ના મુક્ત ચલો પર સંમત હોય, તો તે~$!B$ના
મુક્ત ચલો પર પણ સંમત છે અને અનુમાનાત્મક પૂર્વધારણા~$!B$ને લાગુ પડે છે.

\begin{enumerate}
\tagitem{defNot}{}{%
  \iftag{probNot}{%
    \indcase!{!A}{\lnot !B}{}}{%
    \indcase{!A}{\lnot !B}{જો $\Sat{M}{\indfrm}[s_1]$, તો
      $\Sat/{M}{!B}[s_1]$; તેથી અનુમાનાત્મક પૂર્વધારણાથી
      $\Sat/{M}{!B}[s_2]$, અને પરિણામે $\Sat{M}{\indfrm}[s_2]$.}}}

\tagitem{defAnd}{}{%
  \iftag{probAnd}{%
    \indcase!{!A}{!B \land !C}{}}{%
    \indcase{!A}{!B \land !C}{જો $\Sat{M}{\indfrm}[s_1]$, તો
      $\Sat{M}{!B}[s_1]$ અને $\Sat{M}{!C}[s_1]$; તેથી અનુમાનાત્મક
      પૂર્વધારણાથી $\Sat{M}{!B}[s_2]$ અને $\Sat{M}{!C}[s_2]$.
      પરિણામે $\Sat{M}{\indfrm}[s_2]$.}}}

\tagitem{defOr}{}{%
  \iftag{probOr}{%
    \indcase!{!A}{!B \lor !C}{}}{%
    \indcase{!A}{!B \lor !C}{જો $\Sat{M}{\indfrm}[s_1]$, તો
      $\Sat{M}{!B}[s_1]$ અથવા $\Sat{M}{!C}[s_1]$. અનુમાનાત્મક
      પૂર્વધારણાથી $\Sat{M}{!B}[s_2]$ અથવા $\Sat{M}{!C}[s_2]$;
      તેથી $\Sat{M}{\indfrm}[s_2]$.}}}

\tagitem{defIf}{}{%
  \iftag{probIf}{%
    \indcase!{!A}{!B \lif !C}{}}{%
    \indcase{!A}{!B \lif !C}{જો $\Sat{M}{\indfrm}[s_1]$, તો
    $\Sat/{M}{!B}[s_1]$ અથવા $\Sat{M}{!C}[s_1]$. અનુમાનાત્મક
    પૂર્વધારણાથી $\Sat/{M}{!B}[s_2]$ અથવા $\Sat{M}{!C}[s_2]$;
    તેથી $\Sat{M}{\indfrm}[s_2]$.}}}

\tagitem{defIff}{}{%
  \iftag{probIff}{%
    \indcase!{!A}{!B \liff !C}{}}{%
    \indcase{!A}{!B \liff !C}{જો $\Sat{M}{\indfrm}[s_1]$, તો કાં તો
      $\Sat{M}{!B}[s_1]$ અને $\Sat{M}{!C}[s_1]$, અથવા
      $\Sat/{M}{!B}[s_1]$ અને $\Sat/{M}{!C}[s_1]$. અનુમાનાત્મક
      પૂર્વધારણાથી કાં તો $\Sat{M}{!B}[s_2]$ અને $\Sat{M}{!C}[s_2]$,
      અથવા $\Sat/{M}{!B}[s_2]$ અને $\Sat/{M}{!C}[s_2]$. બંને કિસ્સામાં
      $\Sat{M}{\indfrm}[s_2]$.}}}

\tagitem{defEx}{}{%
  \iftag{probEx}{%
    \indcase!{!A}{\lexists[x][!B]}{}}{%
    \indcase{!A}{\lexists[x][!B]}{જો $\Sat{M}{\indfrm}[s_1]$, તો
      એવો $m \in \Domain{M}$ છે કે
      $\Sat{M}{!B}[\Subst{s_1}{m}{x}]$. હવે
      $s_1'=\Subst{s_1}{m}{x}$ અને $s_2'=\Subst{s_2}{m}{x}$ રાખો.
      ~$!B$ના મુક્ત ચલો $x_1$, \dots, $x_n$ અને~$x$માંના છે.
      $s_1'$ અને~$s_2'$, અનુક્રમે $s_1$ અને~$s_2$નાં $x$-રૂપાંતરો છે
      અને પૂર્વધારણાથી $s_1(x_i)=s_2(x_i)$ હોવાથી
      $s_1'(x_i)=s_2'(x_i)$. $s_1'$ અને~$s_2'$ની વ્યાખ્યાથી
      $s_1'(x)=s_2'(x)=m$. ત્યારે અનુમાનાત્મક પૂર્વધારણા~$!B$ તથા
      $s_1'$,~$s_2'$ને લાગુ પડે છે, તેથી $\Sat{M}{!B}[s_2']$.
      $s_2'=\Subst{s_2}{m}{x}$ હોવાથી એવો $m \in \Domain{M}$ છે કે
      $\Sat{M}{!B}[\Subst{s_2}{m}{x}]$, અને તેથી
      $\Sat{M}{\indfrm}[s_2]$.}}}

\tagitem{defAll}{}{%
  \iftag{probAll}{%
    \indcase!{!A}{\lforall[x][!B]}{}}{%
    \indcase{!A}{\lforall[x][!B]}{જો $\Sat{M}{\indfrm}[s_1]$, તો
      દરેક $m \in \Domain{M}$ માટે
      $\Sat{M}{!B}[\Subst{s_1}{m}{x}]$. આપણે બતાવવું છે કે દરેક
      $m \in \Domain{M}$ માટે $\Sat{M}{!B}[\Subst{s_2}{m}{x}]$ પણ.
      તેથી મનસ્વી $m \in \Domain{M}$ લો અને
      $s_1'=\Subst{s_1}{m}{x}$ તથા $s_2'=\Subst{s_2}{m}{x}$ વિચારો.
      આપણને $\Sat{M}{!B}[s_1']$ મળે છે. ~$!B$ના મુક્ત ચલો
      $x_1$, \dots, $x_n$ અને~$x$માંના છે. $s_1'$ અને~$s_2'$, અનુક્રમે
      $s_1$ અને~$s_2$નાં $x$-રૂપાંતરો છે અને પૂર્વધારણાથી
      $s_1(x_i)=s_2(x_i)$ હોવાથી $s_1'(x_i)=s_2'(x_i)$.
      $s_1'$ અને~$s_2'$ની વ્યાખ્યાથી $s_1'(x)=s_2'(x)=m$. ત્યારે
      અનુમાનાત્મક પૂર્વધારણા~$!B$ તથા $s_1'$,~$s_2'$ને લાગુ પડે છે અને
      આપણને $\Sat{M}{!B}[s_2']$ મળે છે. આ દરેક $m \in \Domain{M}$ને
      લાગુ પડે છે, એટલે બધા $m \in \Domain{M}$ માટે
      $\Sat{M}{!B}[\Subst{s_2}{m}{x}]$; તેથી
      $\Sat{M}{\indfrm}[s_2]$.}}}
\sourcecorrection{OLSEM-010}{મૂળની \texttt{assignments.tex},
પંક્તિઓ~181--182માં $s_1'$ અને~$s_2'$ બંનેને અવ્યાખ્યાયિત~$s$માંથી
બનાવ્યાં છે. આ સાબિતી બે નિયુક્તિઓ~$s_1$ અને~$s_2$ સરખાવે છે;
તેથી અનુક્રમે $\Subst{s_1}{m}{x}$ અને $\Subst{s_2}{m}{x}$
પુનઃસ્થાપિત કર્યાં છે.}
\end{enumerate}
અનુમાનથી આપણને મળે છે કે જ્યારે~$!A$ના મુક્ત !!{variable}s
$x_1$, \dots, $x_n$માંના હોય અને $i=1$, \dots,~$n$ માટે
$s_1(x_i)=s_2(x_i)$ હોય, ત્યારે $\Sat{M}{!A}[s_1]$ ત્યારે અને માત્ર
ત્યારે જ્યારે $\Sat{M}{!A}[s_2]$.
\end{proof}

\begin{probtag}{probNot,probOr,probAnd,probIf,probIff,probEx,probAll}
\olref[fol][syn][ass]{prop:satindep}ની સાબિતી પૂરી કરો.
\end{probtag}

\begin{explain}
!!^{sentence}sમાં મુક્ત ચલો નથી; તેથી કોઈ પણ બે ચલ-નિયુક્તિઓ કોઈ પણ
વાક્યના બધા (શૂન્ય) મુક્ત ચલોને એ જ વસ્તુઓ નિયુક્ત કરે છે. હમણાં
સાબિત કરેલા વિધાનનો અર્થ પછી એ થાય છે કે !!a{sentence} કોઈ સંરચનામાં
ચલ-નિયુક્તિ સાપેક્ષ સંતોષાય છે કે નહિ તે નિયુક્તિથી સંપૂર્ણપણે
સ્વતંત્ર છે. આ હકીકત નોંધી લઈએ. તે આગળ આવતી, ચલ-નિયુક્તિનો ઉલ્લેખ
કર્યા વિનાની, !!a{structure}માં !!a{sentence}ના સંતોષની વ્યાખ્યાને
યોગ્ય ઠેરવે છે.
\end{explain}

\begin{cor}
\ollabel{cor:sat-sentence}
જો $!A$ !!a{sentence} હોય અને~$s$ ચલ-નિયુક્તિ હોય, તો દરેક
ચલ-નિયુક્તિ~$s'$ માટે $\Sat{M}{!A}[s]$ ત્યારે અને માત્ર ત્યારે જ્યારે
$\Sat{M}{!A}[s']$.
\end{cor}

\begin{proof}
  કોઈ પણ ચલ-નિયુક્તિ~$s'$ લો. $!A$ !!a{sentence} હોવાથી તેમાં મુક્ત
  ચલો નથી; તેથી દરેક ચલ-નિયુક્તિ~$s'$,~$!A$ના બધા મુક્ત ચલોને
  તુચ્છપણે એ જ વસ્તુઓ નિયુક્ત કરે છે જે~$s$ કરે છે. તેથી
  \olref{prop:satindep}ની શરત સંતોષાય છે અને આપણને
  $\Sat{M}{!A}[s]$ ત્યારે અને માત્ર ત્યારે જ્યારે
  $\Sat{M}{!A}[s']$ મળે છે.
\end{proof}

\begin{defn}
\ollabel{defn:satisfaction}
જો $!A$ !!a{sentence} હોય, તો આપણે કહીએ છીએ કે !!a{structure}
~$\Struct M$,~$!A$ને \emph{સંતોષે} છે, $\Sat{M}{!A}$, ત્યારે અને માત્ર
ત્યારે જ્યારે બધી ચલ-નિયુક્તિઓ~$s$ માટે $\Sat{M}{!A}[s]$.
\end{defn}

જો $\Sat{M}{!A}$ હોય, તો સાદી રીતે \emph{$!A$,~$\Struct{M}$માં સત્ય
છે} એમ પણ કહીએ છીએ. સંતોષનો ખ્યાલ વ્યક્તિગત !!{sentence}sમાંથી
!!{sentence}sના ગણ સુધી સ્વાભાવિક રીતે વિસ્તરે છે.

\begin{defn}
\ollabel{defn:sat}
જો $\Gamma$ !!{sentence}sનો ગણ હોય, તો આપણે કહીએ છીએ કે
!!a{structure}~$\Struct M$,~$\Gamma$ને \emph{સંતોષે} છે,
$\Sat{M}{\Gamma}$, ત્યારે અને માત્ર ત્યારે જ્યારે બધા
$!A \in \Gamma$ માટે $\Sat{M}{!A}$.
\end{defn}
\sourcecorrection{OLSEM-011}{મૂળની \texttt{assignments.tex},
પંક્તિ~241માં ``વાક્યોનો ગણ~$\Gamma$'' પછી~$\Gamma$ બીજી વાર લખાયું
છે. તે કોઈ નવી શરત કે અભિવ્યક્તિ નથી; અતિરિક્ત પુનરાવર્તન દૂર કર્યું
છે.}

\begin{prop}\ollabel{prop:sentence-sat-true}
  $\Struct{M}$ !!a{structure},~$!A$ !!a{sentence} અને~$s$ ચલ-નિયુક્તિ
  હોય. $\Sat{M}{!A}$ ત્યારે અને માત્ર ત્યારે જ્યારે
  $\Sat{M}{!A}[s]$.
\end{prop}

\begin{proof}
મહાવરો.
\end{proof}

\begin{prob}
\olref[fol][syn][ass]{prop:sentence-sat-true} સાબિત કરો.
\end{prob}

\begin{prop}\ollabel{prop:sat-quant}
  ધારો કે $!A(x)$માં માત્ર~$x$ મુક્ત આવે છે અને $\Struct M$
  !!a{structure} છે. ત્યારે:
  \begin{tagenumerate}{prvEx,prvAll}
    \tagitem{prvEx}{$\Sat{M}{\lexists[x][!A(x)]}$ ત્યારે અને માત્ર
      ત્યારે જ્યારે ઓછામાં ઓછી એક ચલ-નિયુક્તિ~$s$ માટે
      $\Sat{M}{!A(x)}[s]$.}{}
    \tagitem{prvAll}{$\Sat{M}{\lforall[x][!A(x)]}$ ત્યારે અને માત્ર
      ત્યારે જ્યારે બધી ચલ-નિયુક્તિઓ~$s$ માટે
      $\Sat{M}{!A(x)}[s]$.}{}
\end{tagenumerate}
\end{prop}

\begin{proof}
મહાવરો.
\end{proof}

\begin{prob}
\olref[fol][syn][ass]{prop:sat-quant} સાબિત કરો.
\end{prob}

\begin{prob}
\DeclareRobustCommand{\VDash}{\mathrel{||}\joinrel\Relbar}
ધારો કે $\Lang L$ !!{function}s વિનાની ભાષા છે. !!{structure}
~$\Struct M$,~!!a{constant}~$c$ અને $a \in \Domain M$ આપેલાં હોય,
ત્યારે $\Struct M[a/c]$ને~$\Struct M$ જેવી જ !!{structure} તરીકે
વ્યાખ્યાયિત કરો, માત્ર $\Assign{c}{M[a/c]}=a$ રાખો.
!!{sentence}s~$!A$ માટે $\Struct M \VDash !A$ આ રીતે વ્યાખ્યાયિત કરો:
\begin{enumerate}
\tagitem{prvFalse}{%
  \indcase{!A}{\lfalse}{$\Struct{M} \VDash \indfrm$ નહિ.}}{}

\tagitem{prvTrue}{%
  \indcase{!A}{\ltrue}{$\Struct{M} \VDash \indfrm$.}}{}

\item \indcase{!A}{\Atom{R}{d_1, \dots, d_n}}{$\Struct M \VDash \indfrm$
  ત્યારે અને માત્ર ત્યારે જ્યારે
  $\langle \Assign{d_1}{M}, \dots, \Assign{d_n}{M} \rangle \in
  \Assign{R}{M}$.}

\item \indcase{!A}{\eq[d_1][d_2]}{$\Struct M \VDash \indfrm$ ત્યારે
  અને માત્ર ત્યારે જ્યારે $\Assign{d_1}{M}=\Assign{d_2}{M}$.}

\tagitem{prvNot}{%
  \indcase{!A}{\lnot !B}{$\Struct{M} \VDash \indfrm$ ત્યારે અને માત્ર
    ત્યારે જ્યારે $\Struct{M} \VDash {!B}$ નહિ.}}{}

\tagitem{prvAnd}{%
  \indcase{!A}{(!B \land !C)}{$\Struct{M} \VDash \indfrm$ ત્યારે અને
    માત્ર ત્યારે જ્યારે $\Struct{M} \VDash!B$ અને
    $\Struct{M} \VDash !C$.}}{}

\tagitem{prvOr}{%
  \indcase{!A}{(!B \lor !C)}{$\Struct{M} \VDash \indfrm$ ત્યારે અને
    માત્ર ત્યારે જ્યારે $\Struct{M} \VDash !B$ અથવા
    $\Struct{M} \VDash !C$ (અથવા બંને).}}{}

\tagitem{prvIf}{%
  \indcase{!A}{(!B \lif !C)}{$\Struct{M} \VDash \indfrm$ ત્યારે અને
    માત્ર ત્યારે જ્યારે $\Struct {M} \VDash !B$ નહિ અથવા
    $\Struct M \VDash !C$ (અથવા બંને).}}{}

\tagitem{prvIff}{%
  \indcase{!A}{(!B \liff !C)}{$\Struct{M} \VDash {\indfrm}$ ત્યારે
    અને માત્ર ત્યારે જ્યારે $\Struct{M} \VDash {!B}$ અને
    $\Struct{M} \VDash {!C}$ બંને, અથવા $\Struct{M} \VDash {!B}$ પણ
    નહિ અને $\Struct{M} \VDash {!C}$ પણ નહિ.}}{}

\tagitem{prvAll}{%
  \indcase{!A}{\lforall[x][!B]}{$\Struct{M} \VDash {\indfrm}$ ત્યારે
    અને માત્ર ત્યારે જ્યારે બધા $a \in \Domain{M}$ માટે
    $\Struct{M[a/c]} \VDash \Subst{!B}{c}{x}$, જો~$c$,~$!B$માં આવતું
    ન હોય.}}{}

\tagitem{prvEx}{%
  \indcase{!A}{\lexists[x][!B]}{$\Struct{M} \VDash {\indfrm}$ ત્યારે
    અને માત્ર ત્યારે જ્યારે એવું $a \in \Domain M$ હોય કે
    $\Struct{M[a/c]} \VDash \Subst{!B}{c}{x}$, જો~$c$,~$!B$માં આવતું
    ન હોય.}}{}
\end{enumerate}
~$!A$માંના બધા મુક્ત !!{variable}s $x_1$, \dots,~$x_n$,~$!A$માં ન
આવતા અચળસંકેતો $c_1$, \dots,~$c_n$,~$a_1$, \dots,
$a_n \in \Domain M$ અને $s(x_i)=a_i$ હોય.

આ બતાવો:
$\Sat{M}{!A}[s]$ ત્યારે અને માત્ર ત્યારે જ્યારે
$\Struct M[a_1/c_1,\dots,a_n/c_n]
\VDash \Subst{\Subst{!A}{c_1}{x_1}\dots}{c_n}{x_n}$.

(આ પ્રશ્ન બતાવે છે કે ચલ-નિયુક્તિઓ વિનાનો પ્રથમ-ક્રમ તર્કશાસ્ત્રનો
અર્થવિચાર આપી શકાય છે.)
\end{prob}

\begin{prob}
ધારો કે $f$ એ~$!A(x,y)$માં ન આવતો વિધેયસંકેત છે. બતાવો કે એવી
!!a{structure}~$\Struct{M}$ છે કે
$\Sat{M}{\lforall[x][\lexists[y][!A(x,y)]]}$ ત્યારે અને માત્ર ત્યારે
જ્યારે એવી~$\Struct M'$ છે કે
$\Sat{M'}{\lforall[x][!A(x,f(x))]}$.

(આ પ્રશ્ન સ્કોલેમના પ્રમેય તરીકે જાણીતા પરિણામનો એક ખાસ કિસ્સો છે;
$\lforall[x][!A(x,f(x))]$ને
$\lforall[x][\lexists[y][!A(x,y)]]$નું \emph{સ્કોલેમ સામાન્ય સ્વરૂપ}
કહે છે.)
\end{prob}

\end{document}
